%% kapa: exam=1 round=31 year=2020 period=2 form=A subject=appraisal_law subject_ko=감정평가관계법규 seq=34 number=34
%% source: Q-Net 감정평가사 기출문제 (cst00309) artlSeq=5207219 file=extracted/1008/5207219/A형/A형/2020년도 제31회 감정평가사 1차 시험_2교시_A형.pdf
%% conversion: kordoc 4.24.51 -> tools/gen1.py (auto); review=proofread-sonnet; flags=none
\begin{question}{34}
부동산등기법령상 등기신청에 관한 설명으로 옳지 않은 것은?
\begin{choices}
  \item 법인의 합병에 따른 등기는 등기권리자가 단독으로 신청한다.
  \item 신탁재산에 속하는 부동산의 신탁등기는 수탁자가 단독으로 신청한다.
  \item 수용으로 인한 소유권이전등기는 등기권리자가 단독으로 신청할 수 있다.
  \item 채권자는 「민법」 제404조에 따라 채무자를 대위하여 등기를 신청할 수 있다.
  \item 대표자가 있는 종중에 속하는 부동산의 등기는 대표자의 명의로 신청한다.
\end{choices}
\end{question}
\answer{5}
